\documentclass[10pt,a4paper]{amsart}
\usepackage[margin=19mm]{geometry}
\usepackage{amsmath,amssymb,mathtools}
\usepackage[scr=rsfs,cal=euler]{mathalfa}
\usepackage{xfrac}
\usepackage[unicode,hidelinks,bookmarks=false]{hyperref}
\newcommand{\ie}{i.e.\ }
\newcommand{\cf}{cf.\ }
\newcommand{\resp}{respectively\ }
\newcommand{\Subsection}{\subsection}
\providecommand{\nobreakdash}{\nobreak\hbox{-}}
\setlength{\emergencystretch}{2em}
\multlinegap0pt
\usepackage{amssymb}
\usepackage{mathtools}
\usepackage[shortlabels]{enumitem}
\usepackage{tikz}
\newtheorem{lemma}{Lemma}
\newcommand{\E}{\mathbb{E}}
\newcommand{\N}{\mathbb{N}}
\newcommand{\R}{\mathbb{R}}
\newcommand{\Var}{\mathrm{Var}}
\newcommand{\abs}[1]{\left\lvert #1 \right\rvert}
\newcommand{\what}[1]{\widehat{#1}}
\title{Microergodicity implies orthogonality of Mat\'ern fields on bounded domains in $\mathbb{R}^4$}
\author{Natesh S. Pillai}
\date{Source: arXiv:2603.23959v1 (CC BY 4.0)}
\begin{document}
\maketitle

\noindent\emph{Source and licence.} Microergodicity implies orthogonality of Mat\'ern fields on bounded domains in $\mathbb{R}^4$. Version arXiv:2603.23959v1 (CC BY 4.0), distributed by arXiv under the Creative Commons Attribution 4.0 International licence, http://creativecommons.org/licenses/by/4.0/, as recorded in the arXiv licence metadata for this exact version and shown on the abstract page https://arxiv.org/abs/2603.23959v1. Full licence notice: \texttt{licenses/arxiv-2603.23959-CC-BY-4.0.txt}.\par\smallskip

\noindent\textit{Editorial scope.} Counted unit: the article's lemma lem:matdiag (source0.tex lines 748--797). The article's complete original proof of it is delivered with it (source0.tex lines 799--1010, one \texttt{proof} environment of 212 lines). No mathematics is written, repaired or re-derived: the statement and the proof are the article's own lines, in the article's own notation, and no retained line is substituted or re-flowed.  Retained uncounted as the source-local context the proof needs: lines 1409--1431.  Nothing was omitted from any retained span, so the package carries no omission record.  No retained line contains a citation, so the wrapper carries no bibliography and no bibliography entry.  No retained line contains a figure environment, an image-inclusion command or drawing code, so no image is delivered and the package declares no asset dependency.  Editorial formatting only, in addition: an AMS \texttt{amsart} preamble that restates the article's own declarations and macros used by the retained lines, namely the environments \texttt{lemma} on separate counters, together with the standard packages those lines need.  The retained environments print in this document as Lemma 1 in order of appearance, whereas the article prints the same items with the numbering of its own full document.  The complete native source of the article as distributed in the arXiv e-print for this exact version is bundled unchanged as \texttt{sources/197123/source0.tex}.
\begin{lemma}\label{lem:diagonal-asymptotics}
Define
\[
c_\chi:=(2\pi)^{-4}\int_{\R^4}\abs{\what\chi(\eta)}^2\,d\eta = \int_D \chi(t)^2\,dt >0.
\]
Then, as $\abs{k}\to\infty$ with $k\in\N^4$,
\begin{align}
v_j(k)&=c_\chi f_j(k)(1+o(1)),
\label{eq:vi-asymptotic}
\\
v_2(k)-v_1(k)&=c_\chi h(k)(1+o(1)).
\label{eq:diff-asymptotic}
\end{align}
Consequently, if
\begin{equation}\label{eq:delta-def}
\delta_k:=\frac{v_2(k)}{v_1(k)}-1,
\end{equation}
then
\begin{equation}\label{eq:delta-asymptotic}
\delta_k = p(\alpha_1^2-\alpha_2^2)\abs{k}^{-2}+o(\abs{k}^{-2}).
\end{equation}
In particular, $\abs{\delta_k}\asymp \abs{k}^{-2}$ for all sufficiently large $k$.
\end{lemma}

\begin{lemma}[Diagonal likelihood expansion for the localized Mat\'ern coefficients]
\label{lem:matdiag}
Fix \(m>0\) and let \(\alpha_1\neq \alpha_2\). For \(j=1,2\), let \(\mathbf P_j^{\mathrm{diag}}\) denote the diagonal approximation
to the law of \(\{X_k\}_{k\in\Lambda_N}\) under the Mat\'ern model with parameters
\((m,\alpha_j)\), that is, the law under which the variables
\(\{X_k\}_{k\in\Lambda_N}\) are independent, centered, circular complex Gaussian with
\[
\E_j^{\mathrm{diag}} |X_k|^2 = v_j(k), \qquad k\in\Lambda_N.
\]
Then
\begin{equation}\label{eq:matern-diagonal-llr}
\log\frac{d\mathbf P_2^{\mathrm{diag}}}{d\mathbf P_1^{\mathrm{diag}}}
=
S_N-\frac12L_N+R_N
\end{equation}
where
\[
R_N=o_{\mathbf P_j^{\mathrm{diag}}}(L_N), \qquad j=1,2.
\]

We have,
\[
\E_1^{\mathrm{diag}}[S_N]=0,
\qquad
\Var_1^{\mathrm{diag}}(S_N)=L_N,
\]
and
\[
\E_2^{\mathrm{diag}}[S_N]=L_N,
\qquad
\Var_2^{\mathrm{diag}}(S_N)
=
\sum_{k\in\Lambda_N}\delta_k^2(1+\delta_k)^2 \sim L_N.
\]
%Finally,
%\[
%L_N\asymp \sum_{k\in\Lambda_N}|k|^{-4}\asymp \log N,
%\qquad
%\sum_{k\in\Lambda_N}|\delta_k|^3=O(1),
%\]
%and therefore
Moreover
\[
T_N=\frac{S_N}{L_N}\to 0
\quad\text{in \(\mathbf P_1^{\mathrm{diag}}\)-probability},
\qquad
T_N=\frac{S_N}{L_N}\to 1
\quad\text{in \(\mathbf P_2^{\mathrm{diag}}\)-probability}.
\]
\end{lemma}

\begin{proof}
For each \(j=1,2\), the diagonal approximation treats the coefficients
\(\{X_k\}_{k\in\Lambda_N}\) as independent centered circular complex Gaussian
variables with variances \(v_j(k)\). Since the density of a centered circular
complex Gaussian variable with variance \(v\) is
\[
\varphi_v(z)=\frac1{\pi v}\exp\!\left(-\frac{|z|^2}{v}\right), \qquad z\in\mathbb C,
\]
we obtain
\[
\log\frac{d\mathbf P_2^{\mathrm{diag}}}{d\mathbf P_1^{\mathrm{diag}}}
=
\sum_{k\in\Lambda_N}
\left[
-\log\frac{v_2(k)}{v_1(k)}
+
|X_k|^2\left(\frac1{v_1(k)}-\frac1{v_2(k)}\right)
\right].
\]
Using \(v_2(k)=v_1(k)(1+\delta_k)\), this becomes
\[
\log\frac{d\mathbf P_2^{\mathrm{diag}}}{d\mathbf P_1^{\mathrm{diag}}}
=
\sum_{k\in\Lambda_N}
\left[
-\log(1+\delta_k)
+
\frac{|X_k|^2}{v_1(k)}\frac{\delta_k}{1+\delta_k}
\right].
\]
For \(|\delta|\) small,
\[
\log(1+\delta)=\delta-\frac12\delta^2+O(\delta^3),
\qquad
\frac{\delta}{1+\delta}=\delta-\delta^2+O(\delta^3),
\]
hence
\begin{align*}
\log\frac{d\mathbf P_2^{\mathrm{diag}}}{d\mathbf P_1^{\mathrm{diag}}}
&=
\sum_{k\in\Lambda_N}
\delta_k\left(\frac{|X_k|^2}{v_1(k)}-1\right)
-\frac12\sum_{k\in\Lambda_N}\delta_k^2 \\
&\quad
-\sum_{k\in\Lambda_N}\delta_k^2
\left(\frac{|X_k|^2}{v_1(k)}-1\right)
+
E_N,
\end{align*}
where
\[
|E_N|
\le
C\sum_{k\in\Lambda_N}
\left(1+\frac{|X_k|^2}{v_1(k)}\right)|\delta_k|^3.
\]
Thus \eqref{eq:matern-diagonal-llr} holds with
\[
R_N
:=
-\sum_{k\in\Lambda_N}\delta_k^2
\left(\frac{|X_k|^2}{v_1(k)}-1\right)
+
E_N.
\]
We next show that \(R_N=o_{\mathbf P_j^{\mathrm{diag}}}(L_N)\) for \(j=1,2\).
Set
\[
Y_k:=\frac{|X_k|^2}{v_1(k)}-1.
\]
Then
\[
R_N=-\sum_{k\in\Lambda_N}\delta_k^2Y_k+E_N.
\]

Under \(\mathbf P_1^{\mathrm{diag}}\), the variable \(X_k/v_1(k)^{1/2}\) is
standard centered circular complex Gaussian, so
\[
\E_1^{\mathrm{diag}}[Y_k]=0,
\qquad
\Var_1^{\mathrm{diag}}(Y_k)=1.
\]
By independence,
\[
\Var_1^{\mathrm{diag}}\!\left(\sum_{k\in\Lambda_N}\delta_k^2Y_k\right)
=
\sum_{k\in\Lambda_N}\delta_k^4.
\]
Under \(\mathbf P_2^{\mathrm{diag}}\),
\[
\E_2^{\mathrm{diag}}[Y_k]=\delta_k,
\qquad
\Var_2^{\mathrm{diag}}(Y_k)=(1+\delta_k)^2,
\]
and therefore
\[
\Var_2^{\mathrm{diag}}\!\left(\sum_{k\in\Lambda_N}\delta_k^2Y_k\right)
=
\sum_{k\in\Lambda_N}\delta_k^4(1+\delta_k)^2.
\]
By Lemma~\ref{lem:diagonal-asymptotics},
\[
\delta_k=p(\alpha_1^2-\alpha_2^2)|k|^{-2}+o(|k|^{-2}),
\]
so
\[
\sum_{k\in\Lambda_N}\delta_k^4
\asymp
\sum_{k\in\Lambda_N}|k|^{-8}
=O(1),
\qquad
L_N=\sum_{k\in\Lambda_N}\delta_k^2\asymp \log N.
\]
Hence
\[
\sum_{k\in\Lambda_N}\delta_k^2Y_k
=
O_{\mathbf P_j^{\mathrm{diag}}}(1),
\qquad j=1,2.
\]
For the cubic error term, again by \(\delta_k=O(|k|^{-2})\),
\[
\sum_{k\in\Lambda_N}|\delta_k|^3
\asymp
\sum_{k\in\Lambda_N}|k|^{-6}
=
O(1).
\]
Moreover, under each \(\mathbf P_j^{\mathrm{diag}}\), the variables
\(|X_k|^2/v_1(k)\) have uniformly bounded first moments, since
\[
\E_1^{\mathrm{diag}}\!\left[\frac{|X_k|^2}{v_1(k)}\right]=1,
\qquad
\E_2^{\mathrm{diag}}\!\left[\frac{|X_k|^2}{v_1(k)}\right]
=
\frac{v_2(k)}{v_1(k)}=1+\delta_k.
\]
Therefore
\[
\E_j^{\mathrm{diag}}|E_N|
\le
C\sum_{k\in\Lambda_N}
\left(1+\E_j^{\mathrm{diag}}\!\left[\frac{|X_k|^2}{v_1(k)}\right]\right)|\delta_k|^3
=O(1),
\qquad j=1,2,
\]
and hence
\[
E_N=O_{\mathbf P_j^{\mathrm{diag}}}(1),
\qquad j=1,2.
\]
Combining the two bounds gives
\[
R_N=O_{\mathbf P_j^{\mathrm{diag}}}(1),
\qquad j=1,2.
\]
Since \(L_N\asymp \log N\to\infty\), it follows that
\[
R_N=o_{\mathbf P_j^{\mathrm{diag}}}(L_N),
\qquad j=1,2.
\]
It remains to verify the stated moments of \(S_N\). Since
\[
S_N=\sum_{k\in\Lambda_N}\delta_kY_k,
\]
the above calculations give
\[
\E_1^{\mathrm{diag}}[S_N]=0,
\qquad
\Var_1^{\mathrm{diag}}(S_N)=\sum_{k\in\Lambda_N}\delta_k^2=L_N,
\]
and
\[
\E_2^{\mathrm{diag}}[S_N]
=
\sum_{k\in\Lambda_N}\delta_k\,\E_2^{\mathrm{diag}}[Y_k]
=
\sum_{k\in\Lambda_N}\delta_k^2
=
L_N,
\]
while
\[
\Var_2^{\mathrm{diag}}(S_N)
=
\sum_{k\in\Lambda_N}\delta_k^2\Var_2^{\mathrm{diag}}(Y_k)
=
\sum_{k\in\Lambda_N}\delta_k^2(1+\delta_k)^2.
\]

Finally,
\[
\Var_1^{\mathrm{diag}}\!\left(\frac{S_N}{L_N}\right)=\frac1{L_N}\to0,
\]
and
\[
\Var_2^{\mathrm{diag}}\!\left(\frac{S_N}{L_N}\right)
=
\frac{1}{L_N^2}\sum_{k\in\Lambda_N}\delta_k^2(1+\delta_k)^2
\to0,
\]
since \(\sum \delta_k^2(1+\delta_k)^2\asymp L_N\). Hence Chebyshev's inequality
yields
\[
T_N=\frac{S_N}{L_N}\to 0
\quad\text{in \(\mathbf P_1^{\mathrm{diag}}\)-probability},
\qquad
T_N=\frac{S_N}{L_N}\to 1
\quad\text{in \(\mathbf P_2^{\mathrm{diag}}\)-probability}
\]
and the proof is finished.
\end{proof}

\end{document}
